\begin{table}[htbp]
\centering
\caption{Konsekuensi pemilihan ukuran potongan dokumen.}
\label{tab:ukuran-potongan}
\small
\begin{tabular}{@{}p{3.4cm}p{5.4cm}p{5.4cm}@{}}
\toprule
\textbf{Aspek} & \textbf{Potongan berukuran kecil} & \textbf{Potongan berukuran besar} \\
\midrule
Ketajaman pencarian & Lebih tajam karena setiap potongan memuat gagasan yang lebih sempit & Menurun karena satu potongan memuat beberapa gagasan sekaligus \\[2pt]
Keutuhan gagasan & Berisiko memutus satu rekomendasi menjadi beberapa bagian & Lebih terjaga karena kondisi, tindakan, dan peringatan cenderung berada dalam satu potongan \\[2pt]
Derau pada konteks & Rendah karena isi potongan lebih terfokus & Meningkat karena informasi yang tidak relevan turut terbawa \\[2pt]
Jumlah potongan & Banyak & Sedikit \\[2pt]
Biaya pemrosesan & Bertambah pada tahap pengindeksan seiring bertambahnya jumlah potongan & Bertambah pada tahap pembangkitan karena konteks yang dikirimkan lebih panjang \\
\bottomrule
\end{tabular}

\vspace{2pt}
\footnotesize Disarikan dari \textcite{gao2023}.\normalsize
\end{table}
